\documentclass[11pt,a4paper]{article}

% ---------------------------- Packages ------------------------------------
\usepackage[utf8]{inputenc}
\usepackage[T1]{fontenc}
\usepackage{lmodern}
\usepackage[english]{babel}
\usepackage{geometry}
\geometry{margin=2.5cm}

\usepackage{amsmath,amssymb,amsthm}
\usepackage{physics}                 % \dd, \abs, etc.
\usepackage{booktabs}
\usepackage{array}
\usepackage{xcolor}
\usepackage{hyperref}
\usepackage{listings}
\usepackage{tcolorbox}
\tcbuselibrary{breakable,skins}
\usepackage{siunitx}
\usepackage{enumitem}

% ---------------------------- Style ---------------------------------------
\hypersetup{
    colorlinks=true,
    linkcolor=blue!60!black,
    urlcolor=blue!60!black,
    citecolor=blue!60!black
}

% Custom boxes
\newtcolorbox{filebox}[1][]{
    enhanced, breakable,
    colback=gray!5, colframe=gray!50!black,
    fonttitle=\bfseries\small,
    title={Source file},
    boxrule=0.6pt, arc=2pt,
    #1
}

\newtcolorbox{resultbox}[1][]{
    enhanced, breakable,
    colback=blue!4, colframe=blue!50!black,
    fonttitle=\bfseries\small,
    title={Result},
    boxrule=0.6pt, arc=2pt,
    #1
}

\newtcolorbox{warningbox}[1][]{
    enhanced, breakable,
    colback=orange!8, colframe=orange!60!black,
    fonttitle=\bfseries\small,
    title={Convention warning},
    boxrule=0.6pt, arc=2pt,
    #1
}

% Code listing style
\lstset{
    basicstyle=\ttfamily\small,
    breaklines=true,
    frame=single,
    framerule=0.4pt,
    rulecolor=\color{gray!60},
    backgroundcolor=\color{gray!5},
    columns=fullflexible,
    showstringspaces=false
}

% Compact math
\newcommand{\Ef}{E_{\mathrm{F}}}
\newcommand{\Tid}{\text{Ti-}d}
\newcommand{\rhoup}{\rho^{\uparrow}}
\newcommand{\rhodn}{\rho^{\downarrow}}
\newcommand{\Dex}{\Delta_{\text{ex}}}
\newcommand{\Ieff}{I_{\text{eff}}}
\newcommand{\NEf}{N(\Ef)}

% Title
\title{\bfseries
    Stoner Criterion Evaluation\\
    \large An operational guide for (Nb$_{0.8}$Ti$_{0.2}$)$_4$C$_3$
}
\author{}
\date{\today}

% =========================================================================
\begin{document}
\maketitle

\begin{abstract}
\noindent
This document is an operational guide. It shows step by step \emph{how}
the Stoner criterion was evaluated for the 2D MXene
(Nb$_{0.8}$Ti$_{0.2}$)$_4$C$_3$. For each numerical value, we indicate
which Quantum ESPRESSO output file contains it and which formula is
applied. The goal is full reproducibility: anyone who runs the same
calculations should obtain the same numbers (within $\sim 1\%$
numerical noise).
\end{abstract}

% =========================================================================
\section{Introduction}

\subsection{What is the Stoner criterion?}

The Stoner criterion is a band-theory test. It predicts when a metal
spontaneously becomes ferromagnetic. The condition is:
%
\begin{equation}
    \boxed{\,\Ieff \cdot \NEf \;>\; 1\,}
    \label{eq:stoner}
\end{equation}
%
where
\begin{itemize}[itemsep=2pt]
    \item $\NEf$ is the density of states at the Fermi level (per atom,
          per spin) computed in the \emph{non-magnetic} state.
    \item $\Ieff$ is an effective exchange parameter with units of energy.
\end{itemize}

If the product exceeds one, the paramagnetic state is unstable against
spontaneous spin polarization, and the system becomes ferromagnetic.

\subsection{System under study}

\begin{table}[h!]
\centering
\begin{tabular}{l l}
\toprule
Composition           & (Nb$_{0.8}$Ti$_{0.2}$)$_4$C$_3$ \\
Unit cell             & 56 atoms (6~Ti, 26~Nb, 24~C) \\
Dimensionality        & 2D MXene with vacuum \\
Code                  & Quantum ESPRESSO 7.4.1 \\
Functional            & PBE (GGA) \\
Hubbard $U$           & $U(\text{Ti-}3d) = 4.3$~eV, $U(\text{Nb-}4d) = 3.4$~eV \\
\bottomrule
\end{tabular}
\caption{Calculation setup.}
\end{table}

\paragraph{Pseudopotentials.} All pseudopotentials are scalar-relativistic
and taken from the SSSP Efficiency library, version 1.3.0:

\begin{table}[h!]
\centering
\small
\begin{tabular}{l c l l}
\toprule
Element & Mass (a.m.u.) & File & Type \\
\midrule
Ti & 47.867 & \texttt{ti\_pbe\_v1.4.uspp.F.UPF}              & USPP \\
Nb & 92.906 & \texttt{Nb.pbe-spn-kjpaw\_psl.0.3.0.UPF}       & PAW  \\
C  & 12.011 & \texttt{C.pbe-n-kjpaw\_psl.1.0.0.UPF}          & PAW  \\
\bottomrule
\end{tabular}
\caption{Pseudopotentials used in the calculations.}
\end{table}

% =========================================================================
\section{How $\Ieff$ is obtained}

We cannot extract $\Ieff$ directly from a DFT calculation. We use the
standard approximation: in the ferromagnetic state, the exchange
interaction shifts the spin-up and spin-down bands apart by an amount
$\Dex$, and this shift is proportional to the magnetic moment:
%
\begin{equation}
    \Dex \;\approx\; \Ieff \cdot m
    \quad\Longrightarrow\quad
    \boxed{\,\Ieff \;=\; \frac{\Dex}{m_{\Tid}}\,}
    \label{eq:Ieff}
\end{equation}
%
Therefore the criterion in
Eq.~\eqref{eq:stoner} can be rewritten as:
%
\begin{equation}
    \boxed{\,\frac{\Dex}{m_{\Tid}} \cdot \NEf \;>\; 1\,}
    \label{eq:stoner-practical}
\end{equation}

We need \textbf{three numerical inputs}:
%
\begin{enumerate}[label=(\alph*)]
    \item $\NEf$: from the \textbf{non-magnetic} calculation.
    \item $m_{\Tid}$: from the \textbf{ferromagnetic} calculation.
    \item $\Dex$: from the \textbf{ferromagnetic} calculation.
\end{enumerate}

This is why we run \emph{two} independent self-consistent calculations
on the same system.

% =========================================================================
\section{The two calculations we need}

\subsection{Why two calculations?}

The Stoner criterion tests an \emph{instability} of the
\emph{paramagnetic} state. Therefore $\NEf$ must come from the
non-magnetic reference; using $\NEf$ from the spin-polarized result
would underestimate it, because the ferromagnetic state has already
moved states away from $\Ef$.

The local moment $m_{\Tid}$ and the band splitting $\Dex$, on the other
hand, only exist in the ferromagnetic solution.

\subsection{Files produced by each calculation}

Each calculation is a chain of four Quantum ESPRESSO steps:
%
\begin{center}
\texttt{pw.x (scf)} $\to$ \texttt{pw.x (nscf)} $\to$ \texttt{dos.x} $\to$ \texttt{projwfc.x}
\end{center}

\noindent The files we will use are summarized in Table~\ref{tab:files}.
%
\begin{table}[h!]
\small
\centering
\begin{tabular}{l l l}
\toprule
Calculation                    & File                                                  & Provides \\
\midrule
NM (\texttt{nspin=1})          & \texttt{non-magnetic/outputs/bare.opt}                & $\Ef^{\text{NM}}$ \\
NM                             & \texttt{non-magnetic/pdos\_files/mxns.pdos\_atm\#$N$(Ti)\_wfc\#4(d)} & $\rho^{\Tid}(E)$ in NM state \\
FM (\texttt{nspin=2})          & \texttt{magnetic/outputs/bare.opt}                    & $\Ef^{\text{FM}}$, magnetization \\
FM                             & \texttt{magnetic/pdos/mxns.pdos\_atm\#$N$(Ti)\_wfc\#4(d)}            & $\rho^{\uparrow,\Tid}(E)$, $\rho^{\downarrow,\Tid}(E)$ \\
\bottomrule
\end{tabular}
\caption{Files used in the Stoner evaluation. $N$ runs over the six Ti
atoms ($N = 1,\dots,6$).}
\label{tab:files}
\end{table}

% =========================================================================
\section{Step-by-step procedure}

\subsection{Step 1: Extract the two Fermi energies}

The Fermi energy is printed by \texttt{pw.x} at the end of each SCF run.
Search for the line:
%
\begin{lstlisting}
the Fermi energy is X.XXXX ev
\end{lstlisting}

\begin{filebox}
\textbf{NM:} \texttt{non-magnetic/outputs/bare.opt}\\
\textbf{FM:} \texttt{magnetic/outputs/bare.opt}
\end{filebox}

In Python this is one line:
\begin{lstlisting}[language=Python]
matches = re.findall(r"the Fermi energy is\s+([-+.\d]+)\s+ev",
                     scf_out, re.IGNORECASE)
E_F = float(matches[-1])   # take the last (converged) value
\end{lstlisting}

\begin{resultbox}
$\Ef^{\text{NM}} = 5.0476$~eV \\
$\Ef^{\text{FM}} = 5.2374$~eV
\end{resultbox}

\begin{warningbox}
\texttt{dos.x} prints its own Fermi energy in the header of
\texttt{mxns.dos}. Do \textbf{not} use that one — it is recomputed with
the NSCF smearing and may differ by $\sim 10$~meV. Always use the value
from the SCF (\texttt{bare.opt}).
\end{warningbox}

% -------------------------------------------------------------------------
\subsection{Step 2: Compute $\NEf$ from the non-magnetic PDOS}

The non-magnetic calculation produces, for each Ti atom, a file named
\verb|mxns.pdos_atm#N(Ti)_wfc#4(d)| with seven columns. The format is
described in Table~\ref{tab:pdos-NM}.
%
\begin{table}[h!]
\small
\centering
\begin{tabular}{c l}
\toprule
Column & Quantity \\
\midrule
1 & energy $E$ (eV) \\
2 & total Ti-$d$ LDOS (per spin) \\
3 & PDOS $d_{z^2}$ \\
4 & PDOS $d_{xz}$ \\
5 & PDOS $d_{yz}$ \\
6 & PDOS $d_{x^2-y^2}$ \\
7 & PDOS $d_{xy}$ \\
\bottomrule
\end{tabular}
\caption{Column layout of the NM PDOS files (\texttt{nspin=1}, 7
columns).}
\label{tab:pdos-NM}
\end{table}

\paragraph{What to do.} We sum the LDOS over all six Ti atoms, divide by
six (to get a per-atom quantity), and evaluate at $E = \Ef^{\text{NM}}$:
%
\begin{equation}
    \NEf \;=\; \frac{1}{6}\sum_{N=1}^{6}
    \rho^{\text{Ti-}d, \text{NM}}_{N}(\Ef^{\text{NM}}).
    \label{eq:NEf}
\end{equation}

\begin{warningbox}
In \texttt{nspin=1}, the LDOS that Quantum ESPRESSO writes is
\emph{already per spin channel}. The integral $\int_{-\infty}^{\Ef}
\text{LDOS}(E)\,\dd E$ gives $N_{\text{total}}/2$, not $N_{\text{total}}$.
For the Stoner formula, the per-spin convention is exactly what we want
— so we do \textbf{not} multiply by 2.
\end{warningbox}

\paragraph{Pseudo-code.}
\begin{lstlisting}[language=Python]
ldos_d_NM = sum(data_Ti[i][:, 1] for i in range(6))
ldos_d_NM_per_Ti = ldos_d_NM / 6
N_EF = numpy.interp(E_F_NM, E_grid, ldos_d_NM_per_Ti)
\end{lstlisting}

\begin{resultbox}
$\NEf \;=\; 1.1920$~states/eV/Ti/spin
\end{resultbox}

% -------------------------------------------------------------------------
\subsection{Step 3: Compute $m_{\Tid}$ from the ferromagnetic PDOS}

The ferromagnetic PDOS files have 13 columns (spin-resolved): see
Table~\ref{tab:pdos-FM}.
%
\begin{table}[h!]
\small
\centering
\begin{tabular}{c l}
\toprule
Column & Quantity \\
\midrule
1     & energy $E$ (eV) \\
2, 3  & total Ti-$d$ LDOS $\uparrow$, $\downarrow$ \\
4, 5  & PDOS $d_{z^2}$ $\uparrow$, $\downarrow$ \\
6, 7  & PDOS $d_{xz}$  $\uparrow$, $\downarrow$ \\
8, 9  & PDOS $d_{yz}$  $\uparrow$, $\downarrow$ \\
10, 11& PDOS $d_{x^2-y^2}$ $\uparrow$, $\downarrow$ \\
12, 13& PDOS $d_{xy}$ $\uparrow$, $\downarrow$ \\
\bottomrule
\end{tabular}
\caption{Column layout of the FM PDOS files (\texttt{nspin=2}, 13
columns).}
\label{tab:pdos-FM}
\end{table}

\paragraph{What to do.} Sum the LDOS$^{\uparrow}$ and LDOS$^{\downarrow}$
over the six Ti atoms, divide by six, integrate from the lower bound to
$\Ef^{\text{FM}}$, and take the difference:
%
\begin{align}
    N^{\uparrow}_{\Tid}
        &\;=\; \frac{1}{6}\sum_{N=1}^{6}
              \int_{-\infty}^{\Ef^{\text{FM}}}
              \rhoup_{N,\Tid}(E)\,\dd E,
    \label{eq:Nup}\\[4pt]
    N^{\downarrow}_{\Tid}
        &\;=\; \frac{1}{6}\sum_{N=1}^{6}
              \int_{-\infty}^{\Ef^{\text{FM}}}
              \rhodn_{N,\Tid}(E)\,\dd E,
    \label{eq:Ndn}\\[4pt]
    m_{\Tid}
        &\;=\; N^{\uparrow}_{\Tid} - N^{\downarrow}_{\Tid}.
    \label{eq:m}
\end{align}

The integral has units of (states/eV)$\times$eV = electrons. Each
unbalanced electron carries one Bohr magneton, so the difference is
$m_{\Tid}$ in $\mu_{\text{B}}$/Ti.

\paragraph{Pseudo-code.}
\begin{lstlisting}[language=Python]
ldos_up_FM = sum(data_Ti_FM[i][:, 1] for i in range(6)) / 6
ldos_dn_FM = sum(data_Ti_FM[i][:, 2] for i in range(6)) / 6
mask = (E_grid_FM <= E_F_FM)
N_up = numpy.trapezoid(ldos_up_FM[mask], E_grid_FM[mask])
N_dn = numpy.trapezoid(ldos_dn_FM[mask], E_grid_FM[mask])
m_Ti_d = N_up - N_dn
\end{lstlisting}

\begin{resultbox}
$N^{\uparrow}_{\Tid}  = 1.6103$~e/Ti \\
$N^{\downarrow}_{\Tid} = 0.6951$~e/Ti \\
$m_{\Tid} = 0.9152$~$\mu_{\text{B}}$/Ti
\end{resultbox}

\paragraph{Cross-check.} Quantum ESPRESSO also writes Löwdin moments
directly inside the file \texttt{magnetic/outputs/pdos.opt}, in lines
like
\begin{lstlisting}
Atom #  1: ... spin up   = 5.9450, ..., d = 1.6130, dz2= 0.1199, ...
              spin down = 4.9599, ..., d = 0.6996, dz2= 0.0912, ...
\end{lstlisting}
Subtracting $d^{\downarrow}$ from $d^{\uparrow}$ for each Ti and
averaging gives an independent estimate of $m_{\Tid}$ that should agree
with the PDOS integral to within $\sim 1\%$. This is a useful sanity
check.

% -------------------------------------------------------------------------
\subsection{Step 4: Compute the exchange splitting $\Dex$}

We use the \textbf{center-of-mass method}: $\Dex$ is the energy
difference between the centers of mass of the spin-up and spin-down
Ti-$d$ bands.

\paragraph{Formulas.} Let $W = [W_1, W_2]$ be an energy window
(relative to $\Ef^{\text{FM}}$). Define:
%
\begin{align}
    \varepsilon^{\uparrow}_{d}
        &\;=\;
        \frac{\int_{W} E \, \rhoup_{\Tid}(E) \,\dd E}
             {\int_{W} \rhoup_{\Tid}(E) \,\dd E},
    \label{eq:eps-up}\\[4pt]
    \varepsilon^{\downarrow}_{d}
        &\;=\;
        \frac{\int_{W} E \, \rhodn_{\Tid}(E) \,\dd E}
             {\int_{W} \rhodn_{\Tid}(E) \,\dd E},
    \label{eq:eps-dn}\\[4pt]
    \Dex
        &\;=\;
        \varepsilon^{\downarrow}_{d} - \varepsilon^{\uparrow}_{d}.
    \label{eq:Dex}
\end{align}

\paragraph{Window choice.} We use $W = [-8, +10]$~eV (relative to
$\Ef$). This is wide enough to contain the entire Ti-$d$ band on both
spin channels. The robustness to this choice is checked below.

\paragraph{Pseudo-code.}
\begin{lstlisting}[language=Python]
E_rel = E_grid_FM - E_F_FM
mask  = (E_rel >= -8.0) & (E_rel <= 10.0)
eps_up = (numpy.trapezoid(E_rel[mask] * ldos_up_FM[mask], E_rel[mask])
          / numpy.trapezoid(ldos_up_FM[mask], E_rel[mask]))
eps_dn = (numpy.trapezoid(E_rel[mask] * ldos_dn_FM[mask], E_rel[mask])
          / numpy.trapezoid(ldos_dn_FM[mask], E_rel[mask]))
Delta_ex = eps_dn - eps_up
\end{lstlisting}

\begin{resultbox}
$\varepsilon^{\uparrow}_{d} = +0.6173$~eV \\
$\varepsilon^{\downarrow}_{d} = +1.7844$~eV \\
$\Dex = 1.1672$~eV
\end{resultbox}

% -------------------------------------------------------------------------
\subsection{Step 5: Combine the pieces}

We now have everything we need. Following
Eq.~\eqref{eq:Ieff}:
%
\begin{equation}
    \Ieff
    \;=\; \frac{\Dex}{m_{\Tid}}
    \;=\; \frac{1.1672~\text{eV}}{0.9152~\mu_{\text{B}}}
    \;=\; 1.2753~\text{eV}.
\end{equation}

Applying the Stoner criterion of Eq.~\eqref{eq:stoner}:
%
\begin{equation}
    \boxed{\,
    \Ieff \cdot \NEf
    \;=\; 1.2753 \times 1.1920
    \;=\; 1.5201
    \,}
    \label{eq:final}
\end{equation}

\begin{resultbox}
\textbf{Stoner criterion: SATISFIED}\\[2pt]
$\Ieff \cdot \NEf = 1.5201 > 1$, with a margin of $+52.0\%$.\\[2pt]
The non-magnetic state is unstable. The system favours an itinerant
Stoner-type ferromagnetic ordering on the Ti sublattice.
\end{resultbox}

% =========================================================================
\section{Robustness checks}

\subsection{Sensitivity to the window in $\Dex$}

The center-of-mass formula in
Eqs.~\eqref{eq:eps-up}--\eqref{eq:Dex} depends on the chosen energy
window. We have verified that the result is essentially insensitive to
this choice:
%
\begin{table}[h!]
\centering
\begin{tabular}{c c c}
\toprule
Window (eV)      & $\Dex$ (eV)   & Change vs.\ default \\
\midrule
$[-8, +10]$ (default) & 1.1672 & --        \\
$[-7, +8]$            & 1.1672 & 0.00\%    \\
$[-6, +6]$            & 1.1646 & $-0.22\%$ \\
$[-10, +10]$          & 1.1671 & $-0.01\%$ \\
\bottomrule
\end{tabular}
\caption{Exchange splitting under different integration windows.}
\end{table}

\subsection{Cross-check between methods}

The Ti-$d$ moment computed by PDOS integration
($m_{\Tid} = 0.9152$~$\mu_{\text{B}}$/Ti) agrees with the Löwdin moment
in \texttt{magnetic/outputs/pdos.opt} to within $\sim 0.4\%$. This is
the expected level of internal consistency.

% =========================================================================
\section{Physical interpretation}

\begin{itemize}[itemsep=2pt]
    \item The product $\Ieff \cdot \NEf = 1.52$ is well above the
          critical value of 1. The system is unambiguously above
          threshold.
    \item The effective exchange $\Ieff \approx 1.28$~eV is in the range
          typical of early 3$d$ transition metals.
    \item The bandsplitting $\Dex \approx 1.17$~eV is comparable to that
          of bulk iron ($\sim 1.5$~eV) and confirms that the Ti-$d$
          shell is strongly exchange-split.
    \item Because the criterion is satisfied with a $+52\%$ margin, the
          ferromagnetic alignment of the Ti sublattice is robust to
          small variations in calculation parameters (smearing,
          $k$-grid, Hubbard $U$).
\end{itemize}

% =========================================================================
\section{What this criterion does NOT prove}

The Stoner criterion answers \emph{one specific question}: whether the
paramagnetic state is unstable against ferromagnetic alignment on the
relevant orbital. It does \textbf{not} prove the following:

\begin{itemize}[itemsep=2pt]
    \item \textbf{That FM is the ground state.} An antiferromagnetic
          configuration could still be lower in energy. To rule that
          out, we would need to compute candidate AFM configurations
          explicitly.
    \item \textbf{The Curie temperature.} A mean-field Stoner estimate
          systematically overestimates $T_c$. A more reliable value
          requires exchange constants $J_{ij}$ (e.g., from Wannier
          functions) plus a Monte Carlo or magnon spectrum calculation.
    \item \textbf{Half-metallicity.} The system is not half-metallic:
          the spin polarization at $\Ef$ on the total DOS is only
          $\approx 2.4\%$, with both spin channels having finite DOS.
    \item \textbf{Dynamic stability.} The phonon spectrum has not been
          computed in the present work.
\end{itemize}

% =========================================================================
\section{Summary}

\noindent\textbf{Verdict:} Stoner criterion satisfied with a margin of
$+52.0\%$. See the full table of inputs, outputs, and intermediate
quantities below.

\begin{table}[h!]
	\centering
	\small
	\renewcommand{\arraystretch}{1.5}
	\begin{tabular}{l l l}
		\toprule
		Quantity                & Value                & Source \\
		\midrule
		$\Ef^{\text{NM}}$       & 5.0476 eV            & \texttt{non-magnetic/outputs/bare.opt} \\
		$\Ef^{\text{FM}}$       & 5.2374 eV            & \texttt{magnetic/outputs/bare.opt} \\
		$\NEf$                  & 1.1920 states/eV/Ti/spin & NM PDOS Ti-$d$, interpolated at $\Ef^{\text{NM}}$ \\
		$N^{\uparrow}_{\Tid}$   & 1.6103 e/Ti          & FM PDOS, $\int_{-\infty}^{\Ef^{\text{FM}}}\rhoup$ \\
		$N^{\downarrow}_{\Tid}$ & 0.6951 e/Ti          & FM PDOS, $\int_{-\infty}^{\Ef^{\text{FM}}}\rhodn$ \\
		$m_{\Tid}$              & 0.9152 $\mu_{\text{B}}$/Ti & $N^{\uparrow}_{\Tid} - N^{\downarrow}_{\Tid}$ \\
		$\Dex$                  & 1.1672 eV            & center-of-mass, window $[-8, +10]$~eV \\
		$\Ieff$                 & 1.2753 eV            & $\Dex / m_{\Tid}$ \\
		$\Ieff \cdot \NEf$      & \textbf{1.5201}      & Stoner product \\
		\bottomrule
	\end{tabular}
	\caption{Summary of all quantities used and obtained in the Stoner
		evaluation.}
\end{table}

\end{document}
